\documentclass{article}

\usepackage{graphicx}
\usepackage{fancyhdr}
\usepackage[sorting=none]{biblatex}
\usepackage[margin=1in]{geometry}
\usepackage{listings}
\usepackage{tikz}
\usepackage[hidelinks]{hyperref}
\usepackage{subfigure}
\hypersetup{
    colorlinks=true,
    linkcolor=teal,
    filecolor=magenta,      
    urlcolor=teal,
    citecolor = teal
    }
\usepackage{xcolor}
\usepackage{xepersian}
\lstdefinestyle{MATLABStyle}{
    language=Matlab,
    backgroundcolor=\color{white},
    basicstyle=\ttfamily\footnotesize,
    keywordstyle=\color{blue},
    commentstyle=\color{green!50!black},
    stringstyle=\color{red!70!black},
    numbers=left,
    numberstyle=\tiny\color{gray},
    stepnumber=1,
    breaklines=true,
    tabsize=4,
    showstringspaces=false,
    % جلوگیری از تبدیل اعداد به فارسی:
}
\setlength\headheight{28pt} 
\addbibresource{bibliography.bib}
\settextfont[Path={./font/}, Scale=1.3]{IRLotus}
\setlatintextfont[Scale=1]{Times New Roman}
\usepackage{float} 
\renewcommand{\baselinestretch}{1.5}
\pagestyle{fancy}
\fancyhf{}
\rhead{\includegraphics[width=1cm]{img/Logo.png} پاسخ تمرین 2 }
\lhead{\thepage}
\renewcommand{\headrulewidth}{1pt}
\renewcommand{\footrulewidth}{1pt}
\AtBeginDocument{
	\def\chapterautorefname{فصل}%
	\def\sectionautorefname{پاسخ سوال}%
	\def\subsectionautorefname{بخش}%
	\def\subsubsectionautorefname{بخش}%
	\def\equationautorefname{رابطهٔ}%
    \def\lstlistingautorefname{برنامۀ}%
}
\renewcommand{\lstlistingname}{Code}

\definecolor{codegreen}{rgb}{0,0.6,0}
\usepackage{amsmath}
\definecolor{codegray}{rgb}{0.5,0.5,0.5}
\definecolor{codepurple}{rgb}{0.58,0,0.82}
\definecolor{backcolour}{rgb}{0.95,0.95,0.92}

\lstdefinestyle{mystyle}{
	backgroundcolor=\color{backcolour},   
	commentstyle=\color{codegreen},
	keywordstyle=\color{magenta},
	numberstyle=\tiny\color{codegray},
	stringstyle=\color{codepurple},
	basicstyle=\ttfamily\footnotesize,
	breakatwhitespace=false,         
	breaklines=true,                 
	captionpos=b,                    
	keepspaces=true,                 
	numbers=left,                    
	numbersep=5pt,                  
	showspaces=false,                
	showstringspaces=false,
	showtabs=false,                  
	tabsize=2
}

\lstset{style=mystyle}

\begin{document}

\begin{titlepage}
\begin{center}
\defpersianfont\nast[Path={./font/}, Scale=2]{IranNastaliq}
\centerline{{\includegraphics[width=5cm]{img/Logo.png}}}
\centerline{\textcolor[rgb]{0,0,0.5}{\nast \large  دانشگاه صنعتی خواجه نصیرالدین طوسی}}
\centerline{\textcolor[rgb]{0,0,0.5}{\nast \bfseries دانشکدۀ مهندسی برق - گروه مهندسی کنترل}}

\vfill
        
\Huge
\textbf{کنترل خطی}\\
\textbf{پاسخ تمرین 5}\\
        
\vfill
        
\begin{table}[ht]
    \centering
    \huge
    \begin{tabular}{|c|c|}
    \hline
    نام و نام خانوادگی & محمد امین حری فراهانی\\
    \hline
    شمارۀ دانشجویی & 40216674\\
    \hline
    تاریخ & آذر 1404\\
    \hline
    \end{tabular}
\end{table}
\end{center}
\end{titlepage}


\tableofcontents \clearpage
\listoffigures \clearpage
\clearpage
 \clearpage
\newpage



\section{سوال 1}\label{Section1}
با توجه به شکل میتوان فهمید که :

\begin{figure}[h!]
\centering
\begin{minipage}{0.45\textwidth}
    \includegraphics[width=\textwidth]{img/oo.png}
\end{minipage}
\hfill
\begin{minipage}{0.5\textwidth}


\[
\text{(بالا زدگی\%)} \qquad
\scalebox{1.05}{$e^{-\dfrac{\pi \zeta}{\sqrt{1-\zeta^{2}}}}$}
= 0.434
\]

\[
\text{(زمان نشست با معیار 2\%)} \qquad
\scalebox{1.1}{$\displaystyle \frac{4}{\zeta \omega_n}$}
= 1.3
\]


\end{minipage}


\begin{tikzpicture}[overlay, remember picture]

% فلش نارنجی خمیده
\draw[->, orange, thick, line width=2pt]
  (2.2,4.8) .. controls (0.5,6) .. (-2,4);

% فلش بنفش خمیده
\draw[->, violet, thick, line width=2pt]
  (4.5,2.3) .. controls (2,0.5) .. (-1.8,1.5);

\end{tikzpicture}


\end{figure}

خب با توجه به رابطهٔ (1) و چون تنها \(\zeta\) در آن مجهول است، می‌توان مقدار \(\zeta\) را محاسبه کرد که مقدار تقریبی آن برابر است با :

{\Large
\[
\boxed{\zeta \approx 0.256}
\]
}
حال می‌توانیم مقدار \(\zeta\) را در رابطهٔ (2) (زمان نشست) قرار دهیم تا \(\omega_n\) به‌دست آید.
 با قرار دادن مقدار \(\zeta \approx 0.256\) در رابطهٔ (2) داریم:

\[
\frac{4}{\zeta \omega_n} = 1.3
\]

\[
\frac{4}{0.256\, \omega_n} = 1.3
\]
{\Large
\[
\boxed{\omega_n \approx 12.02}
\]
}
{\Large
\[
T(s) = \frac{k \omega_n^{2}}{s^{2} + 2\zeta \omega_n s + \omega_n^{2}} = \frac{144.48k}{s^{2}+6.15s+144.48}
\]
}
حال می‌توانیم \(k\) را از مقدار پیک محاسبه کنیم. اگر \(\zeta\) را در رابطهٔ (1) قرار دهیم، داریم:
{\Large
\[
1 + e^{-\frac{\pi (0.256)}{\sqrt{1 - (0.256)^{2}}}} \approx 1 + 0.435 = 1.435
\]
}

با توجه به اینکه مقدار پیک در شکل \(  2.27\) است، مقدار \(k\) برابر است با:

\[
k = \frac{2.27}{1.435} \approx 1.58
\]
\vspace*{2cm} 
پس در نتیجه \(T(s)\) می‌شود:
{\Large
\[
T(s) = \frac{228.27}{s^{2}+6.15s+144.48}
\]
}

میتوانیم در متلب پاسخ پله این تابع تبدیل را بدست بیاریم و با نتیجه مقایسه کنیم :

\vspace*{1cm} 

{\Large
\textbf{کد متلب:}
}
\begin{LTR}
\begin{lstlisting}[style=MATLABStyle]
clc
clear

zeta = 0.256;
wn  = 12.02;
k   = 1.58;
num = k * wn^2;
den = [1  2*zeta*wn  wn^2];

sys = tf(num, den);

figure;
step(sys)
grid on
title('Step Response of T(s)')
\end{lstlisting}
\end{LTR}

\begin{figure}[H]
  \centering
  \includegraphics[width=0.6\textwidth]{img/pp.png}
  \caption{خروجی متلب}
  \label{fig:block-diagram}
\end{figure}


حالا می‌توانیم \( G(s) \) را حساب کنیم. داریم:
\[
T(s) = \frac{G(s)}{G(s) + 1}  \Rightarrow  G(s) = \frac{T(s)}{1 - T(s)}
\]

\[
T(s)=\frac{k\omega_n^2}{s^2 + 2\zeta\omega_n s + \omega_n^2}, \qquad H(s)=1
\]

\[
G(s)=\frac{T(s)}{1-T(s)}
=\frac{\dfrac{k\omega_n^2}{s^2 + 2\zeta\omega_n s + \omega_n^2}}
{1-\dfrac{k\omega_n^2}{s^2 + 2\zeta\omega_n s + \omega_n^2}}
=\frac{k\omega_n^2}{\,s^2 + 2\zeta\omega_n s + \omega_n^2 - k\omega_n^2\,}.
\]

\text{با جایگذاری مقادیر}
\[
\zeta = 0.256,\quad \omega_n = 12.02,\quad k = 1.58
\]

\[
\omega_n^2 = (12.02)^2 = 144.4804,
\qquad k\omega_n^2 = 1.58\times144.4804 = 228.279032,
\]

\[
2\zeta\omega_n = 2\times0.256\times12.02 = 6.15424,
\qquad \omega_n^2 - k\omega_n^2 = 144.4804 - 228.279032 = -83.798632.
\]

\[
\boxed{\,G(s)=\dfrac{228.279032}{s^{2} + 6.15424\,s - 83.798632}\,}
\]




\section{سوال 2}

\subsection{استفاده از فرمول میسون}



\begin{figure}[h!]
\centering
\begin{minipage}{0.45\textwidth}
    \includegraphics[width=\textwidth]{img/99.png}
\end{minipage}
\hfill
\begin{minipage}{0.5\textwidth}


\[
T(s) = \frac{KG(s)}{KG(s) + 1} 
\]
\end{minipage}
\end{figure}

\[
\Rightarrow T(s)=\frac{k}{s^2 + 4s +3+k}
\]

\[
2\zeta\omega_n = 4 \quad  \zeta = \frac{\sqrt{ 2}} {2}  \Rightarrow \omega_n =2\sqrt{2}
\]

پس باید :
\[
(3 + k) = \omega_{n}^{2} = 8
\]

\[
k = 5
\]



\section{سوال 3}


\begin{figure}[H]
  \centering
  \includegraphics[width=0.9\textwidth]{img/00.png}
  \caption{ بلوک دیاگرام سیستم کنترل وضعیت موتور DC}
  \label{fig:block-diagram}
\end{figure}
\subsection{تابع تبدیل حلقه باز و بسته}

\[
\text{(حلقه بسته)} \qquad \frac{\Theta(s)}{R(s)}= \frac{K_a\frac{K_m}{R_a}\frac{1}{Js + f}\frac{1}{s}}{1 - ( \frac{K_m}{R_a} .\frac{-1}{Js + f}.k_b + K_a . \frac{K_m}{R_a}.\frac{-1}{Js + f}.\frac{1}{s}.k_t )} =\frac{0.02}{s^{2}+0.4s+0.02}
\]

\[
\text{(حلقه باز)} \qquad \frac{\Theta(s)}{R(s)}= \frac{K_a\frac{K_m}{R_a}\frac{1}{Js + f}\frac{1}{s}}{1 - ( \frac{K_m}{R_a} .\frac{-1}{Js + f}.k_b)}  = \frac{\Theta(s)}{R(s)}=\dfrac{0.02}{s(s+0.4)}
\]

\subsection{مدل شبیه سازی شده}

\begin{figure}[H]
  \centering
  \includegraphics[width=0.9\textwidth]{img/09.pdf}
  \label{fig:block-diagram}
\end{figure}


\begin{figure}[H]
  \centering
  \includegraphics[width=0.95\textwidth]{img/com.png}
  \label{fig:block-diagram}
\end{figure}


مشاهده میشود که نتایج متلب با حل دستی یکسان است فقط در متلب صورت و مخرج هر دو تابع تبدیل در 50 ضرب شده 

\subsection{پاسخ پله }

\begin{figure}[h!]
    \centering
    % تصویر اول
    \begin{minipage}[b]{0.45\textwidth}
        \centering
        \includegraphics[width=\textwidth]{img/111.png} % نام فایل تصویر اول
        \caption{فرامیرایی}
        \label{fig:image1}
    \end{minipage}
    \hfill
    % تصویر دوم
    \begin{minipage}[b]{0.45\textwidth}
        \centering
        \includegraphics[width=\textwidth]{img/222.png} % نام فایل تصویر دوم
        \caption{میرایی مرزی}
        \label{fig:image2}
    \end{minipage}
\end{figure}



\begin{figure}[h!]
    \centering
    % تصویر اول
    \begin{minipage}[b]{0.45\textwidth}
        \centering
        \includegraphics[width=\textwidth]{img/333.png} % نام فایل تصویر اول
        \caption{  ناپایدار}
        \label{fig:image1}
    \end{minipage}
    \hfill
    % تصویر دوم
    \begin{minipage}[b]{0.45\textwidth}
        \centering
        \includegraphics[width=\textwidth]{img/444.png} % نام فایل تصویر دوم
         \caption{  ناپایدار}
        \label{fig:image2}
    \end{minipage}
\end{figure}

\subsection*{ج)}
\begin{latin}	
\[
L(s)= \dfrac{0.02}{s^{2}+0.4s} 	
\]
\end{latin}

\begin{latin}
	\[
	k_{v}=\lim_{s\to0}sL(s) = \lim_{s\to0}\frac{0.02}{s^{2}+0.4s} = \frac{1}{20}  \quad \Rightarrow \quad e_{ss}=\frac{1}{K_{v}}=20
	\]
\end{latin}
 خروجی با تابع شیب حدود 20 واحد فاصله دارد.\\

\newpage
\begin{figure}[h]
	\centering
	\includegraphics[scale=0.4]{img/fig1.jpg}
\end{figure}

\subsection*{د)}
فراجهش برابر است با 
$M_{p}=e^{-\dfrac{\zeta\pi}{\sqrt{1-\zeta^{2}}}}$
\\
واضح است که مقدار فراجهش تنها به 
$\zeta$
 وابسته است . 
\\
بار دیگر تابع تبدیل کل سیستم را بر حسب پارامتر 
$K_{a}$
مینویسیم.
\begin{latin}
	$\dfrac{0.8 K_{a}}{s^{2}+0.4s+0.4K_{a}}$
\end{latin}
\begin{latin}
	\[
	\omega_{n}^{^{2}}=0.4K_{a}
	\qquad
	2\omega_{n}\zeta=0.4
	\qquad\Rightarrow \zeta =\frac{0.2}{\sqrt{0.4K_{a}}}
	\]
\end{latin}
\begin{latin}
\[
M_{p}<5 \quad \rightarrow \quad e^{\dfrac{-\zeta\pi}{\sqrt{1-\zeta^{2}}}}<0.05 \quad\rightarrow\quad\dfrac{-\zeta\pi}{\sqrt{1-\zeta^{2}}}<\ln0.05
\quad \Rightarrow \zeta > 0.68
\]
\end{latin}
\begin{latin}
	\[
		\Rightarrow\qquad
	\dfrac{0.2}{\sqrt{0.4K_{a}}} >0.68
	\Rightarrow K_{a}<0.21
	\]
\end{latin}
\newpage
پاسخ پله را در متلب به ازای مقادیر مختلف ka مشاهده میکنید.
\\
\begin{figure}[h]
	\centering
	\includegraphics[scale=0.25]{img/figd.jpg}
\end{figure}
همانطور که محاسبه کردیم به ازای $k_{a}$
بزرگتر از \lr{0.21} overshoot بزرگتر از ۵ درصد خواهد بود.





\section{ضمیمه}
\href{https://github.com/aminism00/LinearControl2025/tree/main/HWs/HW3}{ ادرس گیت هاب  (\lr{GitHub})}.



\end{document}
